\documentclass[10pt,a4paper]{amsart}
\usepackage{amsmath,amssymb,amsthm,graphicx}
\usepackage{verbatim}
\newenvironment{reference}{\comment}{\endcomment}
%\newenvironment{reference}{}{}
\newenvironment{slogan}{\comment}{\endcomment}
\newenvironment{history}{\comment}{\endcomment}
% For commutative diagrams we use Xy-pic
\usepackage[all]{xy}

% We use 2cell for 2-commutative diagrams.
\xyoption{2cell}
% Theorem environments.
%
\theoremstyle{plain}
\newtheorem{theorem}[subsection]{Theorem}
\newtheorem{proposition}[subsection]{Proposition}
\newtheorem{lemma}[subsection]{Lemma}

\theoremstyle{definition}
\newtheorem{definition}[subsection]{Definition}
\newtheorem{example}[subsection]{Example}
\newtheorem{exercise}[subsection]{Exercise}
\newtheorem{situation}[subsection]{Situation}

\theoremstyle{remark}
\newtheorem{remark}[subsection]{Remark}
\newtheorem{remarks}[subsection]{Remarks}

\numberwithin{equation}{subsection}

% Macros
%
\def\lim{\mathop{\mathrm{lim}}\nolimits}
\def\colim{\mathop{\mathrm{colim}}\nolimits}
\def\Spec{\mathop{\mathrm{Spec}}}
\def\Hom{\mathop{\mathrm{Hom}}\nolimits}
\def\Ext{\mathop{\mathrm{Ext}}\nolimits}
\def\SheafHom{\mathop{\mathcal{H}\!\mathit{om}}\nolimits}
\def\SheafExt{\mathop{\mathcal{E}\!\mathit{xt}}\nolimits}
\def\Sch{\mathit{Sch}}
\def\Mor{\mathop{\mathrm{Mor}}\nolimits}
\def\Ob{\mathop{\mathrm{Ob}}\nolimits}
\def\Sh{\mathop{\mathit{Sh}}\nolimits}
\def\NL{\mathop{N\!L}\nolimits}
\def\CH{\mathop{\mathrm{CH}}\nolimits}
\def\proetale{{pro\text{-}\acute{e}tale}}
\def\etale{{\acute{e}tale}}
\def\QCoh{\mathit{QCoh}}
\def\Ker{\mathop{\mathrm{Ker}}}
\def\Im{\mathop{\mathrm{Im}}}
\def\Coker{\mathop{\mathrm{Coker}}}
\def\Coim{\mathop{\mathrm{Coim}}}

% Boxtimes
%
\DeclareMathSymbol{\boxtimes}{\mathbin}{AMSa}{"02}

%
% Macros for moduli stacks/spaces
%
\def\QCohstack{\mathcal{QC}\!\mathit{oh}}
\def\Cohstack{\mathcal{C}\!\mathit{oh}}
\def\Spacesstack{\mathcal{S}\!\mathit{paces}}
\def\Quotfunctor{\mathrm{Quot}}
\def\Hilbfunctor{\mathrm{Hilb}}
\def\Curvesstack{\mathcal{C}\!\mathit{urves}}
\def\Polarizedstack{\mathcal{P}\!\mathit{olarized}}
\def\Complexesstack{\mathcal{C}\!\mathit{omplexes}}
% \Pic is the operator that assigns to X its picard group, usage \Pic(X)
% \Picardstack_{X/B} denotes the Picard stack of X over B
% \Picardfunctor_{X/B} denotes the Picard functor of X over B
\def\Pic{\mathop{\mathrm{Pic}}\nolimits}
\def\Picardstack{\mathcal{P}\!\mathit{ic}}
\def\Picardfunctor{\mathrm{Pic}}
\def\Deformationcategory{\mathcal{D}\!\mathit{ef}}
\usepackage{cleveref}
\crefname{originallemma}{original lemma}{original lemmas}
\Crefname{originallemma}{Original lemma}{Original lemmas}
\makeatletter
\newlabel{lemma-gerbe-decent}{{lemma-gerbe-decent}{source}{}{sourcepointer-1}{}}
\newlabel{lemma-gerbe-decent@cref}{{[originallemma][0][]lemma-gerbe-decent}{[1][0][]source}{}{}{}}
\makeatother
\crefname{originallemma}{original lemma}{original lemmas}
\Crefname{originallemma}{Original lemma}{Original lemmas}
\makeatletter
\newlabel{lemma-gerbe-fppf}{{lemma-gerbe-fppf}{source}{}{sourcepointer-2}{}}
\newlabel{lemma-gerbe-fppf@cref}{{[originallemma][0][]lemma-gerbe-fppf}{[1][0][]source}{}{}{}}
\makeatother
\crefname{originallemma}{original lemma}{original lemmas}
\Crefname{originallemma}{Original lemma}{Original lemmas}
\makeatletter
\newlabel{lemma-representable-decent}{{lemma-representable-decent}{source}{}{sourcepointer-3}{}}
\newlabel{lemma-representable-decent@cref}{{[originallemma][0][]lemma-representable-decent}{[1][0][]source}{}{}{}}
\makeatother
\crefname{originalproposition}{original proposition}{original propositions}
\Crefname{originalproposition}{Original proposition}{Original propositions}
\makeatletter
\newlabel{proposition-decent-sober}{{proposition-decent-sober}{source}{}{sourcepointer-4}{}}
\newlabel{proposition-decent-sober@cref}{{[originalproposition][0][]proposition-decent-sober}{[1][0][]source}{}{}{}}
\makeatother
\crefname{originalproposition}{original proposition}{original propositions}
\Crefname{originalproposition}{Original proposition}{Original propositions}
\makeatletter
\newlabel{proposition-open-stratum}{{proposition-open-stratum}{source}{}{sourcepointer-5}{}}
\newlabel{proposition-open-stratum@cref}{{[originalproposition][0][]proposition-open-stratum}{[1][0][]source}{}{}{}}
\makeatother
\crefname{originalreference}{original reference}{original references}
\Crefname{originalreference}{Original reference}{Original references}
\makeatletter
\newlabel{spaces-over-fields-definition-integral-algebraic-space}{{spaces-over-fields-definition-integral-algebraic-space}{source}{}{sourcepointer-6}{}}
\newlabel{spaces-over-fields-definition-integral-algebraic-space@cref}{{[originalreference][0][]spaces-over-fields-definition-integral-algebraic-space}{[1][0][]source}{}{}{}}
\makeatother
\crefname{originalreference}{original reference}{original references}
\Crefname{originalreference}{Original reference}{Original references}
\makeatletter
\newlabel{spaces-over-fields-section-integral-spaces}{{spaces-over-fields-section-integral-spaces}{source}{}{sourcepointer-7}{}}
\newlabel{spaces-over-fields-section-integral-spaces@cref}{{[originalreference][0][]spaces-over-fields-section-integral-spaces}{[1][0][]source}{}{}{}}
\makeatother
\crefname{originalreference}{original reference}{original references}
\Crefname{originalreference}{Original reference}{Original references}
\makeatletter
\newlabel{stacks-properties-lemma-map-into-reduction}{{stacks-properties-lemma-map-into-reduction}{source}{}{sourcepointer-8}{}}
\newlabel{stacks-properties-lemma-map-into-reduction@cref}{{[originalreference][0][]stacks-properties-lemma-map-into-reduction}{[1][0][]source}{}{}{}}
\makeatother
\title{An integral stack has a sober irreducible space with a unique generic point and a dense open gerbe over an integral scheme}
\author{The Stacks Project Authors}
\date{stacks-project:stacks-morphisms.tex:local-snapshot-20261012; GNU FDL 1.2 or later}
\begin{document}
\maketitle
\noindent\textit{Editorial note.} All delivered mathematical text is the author's own, copied from the pinned Stacks Project chapter source sources/stacks-algebraic-stacks-morphisms/pinned-original-source.tex; no statement or proof step was rewritten, reordered or removed.\par\medskip
\noindent\textit{Editorial note.} Preamble provenance: the Stacks Project keeps one single shared preamble for every chapter; this card takes it directly from the staged sources/stacks-algebraic-stacks-morphisms/preamble.tex as three verbatim spans (lines 8-12, 14-18 and 154-217), so no derived file was written and nothing was added or redefined. The wrapper class and the base packages are supplied by the assembler.\par\medskip
\noindent\textit{Editorial note.} Bibliography provenance: neither the counted statement nor its proof contains a citation (this lane rescans the whole span before writing the card), so the card carries no bibliography and no entry was reconstructed.\par\medskip
\noindent\textit{Editorial note.} Selection provenance: this card was opened under the file-based control contract (CONTROL.md), with the mechanical pre-card gate run on its statement and proof interval before the card was written.\par\medskip
\noindent\textit{Allowed standard prerequisites (editorial).}\begin{itemize}
\item The sobriety of the space of a decent stack with quasi-compact inertia. Invoked at the start of the counted proof (source lines 10660-10664). Gives sobriety and hence a unique generic point.
\item The existence of a dense open substack which is a gerbe over an algebraic space. Invoked inside the counted proof (source lines 10665-10668). Produces the open gerbe of the second assertion.
\item The decency of that base algebraic space and the reduction of the base to an integral scheme. Invoked inside the counted proof (source lines 10669-10683) as cross-chapter pointers. Upgrades the algebraic space to an integral scheme after shrinking.
\end{itemize}\medskip
\ifdefined\hypertarget\hypertarget{sourcepointer-1}{}\fi
\noindent\textit{Source reference.} Original lemma lemma-gerbe-decent: sources/stacks-algebraic-stacks-morphisms/pinned-original-source.tex, source lines 10493--10493. Navigation pointer only: the counted text refers to the Stacks item lemma-gerbe-decent, which is an existing labeled result of this same chapter and is not part of this card. The pointer carries the target's own label as its number and its definition line as its span; its content is not reproduced here.\par\medskip
\ifdefined\hypertarget\hypertarget{sourcepointer-2}{}\fi
\noindent\textit{Source reference.} Original lemma lemma-gerbe-fppf: sources/stacks-algebraic-stacks-morphisms/pinned-original-source.tex, source lines 6287--6287. Navigation pointer only: the counted text refers to the Stacks item lemma-gerbe-fppf, which is an existing labeled result of this same chapter and is not part of this card. The pointer carries the target's own label as its number and its definition line as its span; its content is not reproduced here.\par\medskip
\ifdefined\hypertarget\hypertarget{sourcepointer-3}{}\fi
\noindent\textit{Source reference.} Original lemma lemma-representable-decent: sources/stacks-algebraic-stacks-morphisms/pinned-original-source.tex, source lines 10425--10425. Navigation pointer only: the counted text refers to the Stacks item lemma-representable-decent, which is an existing labeled result of this same chapter and is not part of this card. The pointer carries the target's own label as its number and its definition line as its span; its content is not reproduced here.\par\medskip
\ifdefined\hypertarget\hypertarget{sourcepointer-4}{}\fi
\noindent\textit{Source reference.} Original proposition proposition-decent-sober: sources/stacks-algebraic-stacks-morphisms/pinned-original-source.tex, source lines 10575--10575. Navigation pointer only: the counted text refers to the Stacks item proposition-decent-sober, which is an existing labeled result of this same chapter and is not part of this card. The pointer carries the target's own label as its number and its definition line as its span; its content is not reproduced here.\par\medskip
\ifdefined\hypertarget\hypertarget{sourcepointer-5}{}\fi
\noindent\textit{Source reference.} Original proposition proposition-open-stratum: sources/stacks-algebraic-stacks-morphisms/pinned-original-source.tex, source lines 6746--6746. Navigation pointer only: the counted text refers to the Stacks item proposition-open-stratum, which is an existing labeled result of this same chapter and is not part of this card. The pointer carries the target's own label as its number and its definition line as its span; its content is not reproduced here.\par\medskip
\ifdefined\hypertarget\hypertarget{sourcepointer-6}{}\fi
\noindent\textit{Source reference.} Original reference spaces-over-fields-definition-integral-algebraic-space: sources/stacks-algebraic-stacks-morphisms/pinned-original-source.tex, source lines 10679--10679. Navigation pointer only: the counted text refers to the Stacks item spaces-over-fields-definition-integral-algebraic-space, which belongs to another chapter of the project and is not part of this card. The pointer carries the target's own label as its number because this chapter's staged snapshot does not contain the other chapter's tag table; its content is not reproduced here.\par\medskip
\ifdefined\hypertarget\hypertarget{sourcepointer-7}{}\fi
\noindent\textit{Source reference.} Original reference spaces-over-fields-section-integral-spaces: sources/stacks-algebraic-stacks-morphisms/pinned-original-source.tex, source lines 10630--10630. Navigation pointer only: the counted text refers to the Stacks item spaces-over-fields-section-integral-spaces, which belongs to another chapter of the project and is not part of this card. The pointer carries the target's own label as its number because this chapter's staged snapshot does not contain the other chapter's tag table; its content is not reproduced here.\par\medskip
\ifdefined\hypertarget\hypertarget{sourcepointer-8}{}\fi
\noindent\textit{Source reference.} Original reference stacks-properties-lemma-map-into-reduction: sources/stacks-algebraic-stacks-morphisms/pinned-original-source.tex, source lines 10673--10673. Navigation pointer only: the counted text refers to the Stacks item stacks-properties-lemma-map-into-reduction, which belongs to another chapter of the project and is not part of this card. The pointer carries the target's own label as its number because this chapter's staged snapshot does not contain the other chapter's tag table; its content is not reproduced here.\par\medskip
\begin{lemma}
\label{lemma-structure-integral-stack}
Let $\mathcal{X}$ be an integral algebraic stack. Then
\begin{enumerate}
\item $|\mathcal{X}|$ is sober, irreducible, and has a unique generic point,
\item there exists an open substack $\mathcal{U} \subset \mathcal{X}$ 
which is a gerbe over an integral scheme $U$.
\end{enumerate}
\end{lemma}

\begin{proof}
Proposition \ref{proposition-decent-sober} tells us that $|\mathcal{X}|$
is sober. Of course it is also irreducible and hence has a unique generic
point $x$ (by the definition of sobriety).
Proposition \ref{proposition-open-stratum} shows
the existence of a dense open $\mathcal{U} \subset \mathcal{X}$
which is a gerbe over an algebraic space $U$.
Then $U$ is a decent algebraic space by
Lemma \ref{lemma-gerbe-decent} (and the fact that $\mathcal{U}$
is decent by Lemma \ref{lemma-representable-decent}).
Since $|U| = |\mathcal{U}|$ we see that $|U|$ is irreducible.
Finally, since $\mathcal{U}$ is reduced the morphism
$\mathcal{U} \to U$ factors through $U_{red}$, see
Properties of Stacks, Lemma \ref{stacks-properties-lemma-map-into-reduction}.
Now since $\mathcal{U} \to U$ is flat, locally of finite presentation,
and surjective (Lemma \ref{lemma-gerbe-fppf}), this implies that
$U = U_{red}$, i.e., $U$ is reduced (small detail omitted).
It follows that $U$ is an integral algebraic space, see
Spaces over Fields, Definition
\ref{spaces-over-fields-definition-integral-algebraic-space}.
Then finally, we may replace $U$ (and correspondingly $\mathcal{U}$)
by an open subspace and assume that $U$ is an integral scheme, see
discussion in  Spaces over Fields, Section
\ref{spaces-over-fields-section-integral-spaces}.
\end{proof}

\end{document}
